% ===================================================================
%  CUADRO N.º 16 — Los Grandes Almacenes. — El Bazar. — Los Juguetes.
%  Traduction 2026 d'apres texto/io, controlee sur texto/fr et
%  texto/en. Consigne : voir texto/es/10-cuadro-01.tex.
%
%  Le tableau d'astronomie (bloc 15-1) porte des renvois a LETTRES,
%  (a) a (f), graves sur le tableau lui-meme : ils se rangent sous le
%  numero 85, comme le dit gravuri/literi.json. L'espagnol les compose
%  avec les deux parentheses, comme l'anglais.
% ===================================================================

%%K t16-apar-1 apar
\VUsaut{4.0mm}
\VUtitre{15.0pt}{60}{CUADRO N.º 16}
\VUsaut{2.0mm}
\VUpk{13.2pt}{40}{Los Grandes Almacenes. --- El Bazar.}
\VUpk{13.2pt}{40}{Los Juguetes.}
\VUsaut{3.4mm}

%%K t16-01-1 p
1. --- Se acerca el año nuevo. Los grandes almacenes han preparado sus
exposiciones de \VUgras{juguetes} \textsuperscript{(1)} y regalos de año nuevo.

%%K t16-01-2 p
Pedí a mi abuelo que me llevara al bazar a ver esa hermosa exposición, y
aceptó.

%%K t16-02-1 p
2. --- Primero nos paramos delante de unas hermosas panoplias de armas con
un \VUgras{fusil}, una \VUgras{gorra de plato}, un \VUgras{sable}, un
\VUgras{tahalí} y un par de \VUgras{charreteras}, o bien un \VUgras{casco},
una \VUgras{coraza} \textsuperscript{(3)} y una \VUgras{pistola} \textsuperscript{(4)}. Todo eso
me interesaba mucho más que las \VUgras{linternas mágicas} \textsuperscript{(2)}.

%%K t16-tit-1 sub
\VUsaut{3.0mm}
\VUpk{10.2pt}{40}{Hermosos espectáculos.}
\VUsaut{2.6mm}

%%K t16-03-1 p
3. --- En la misma sección había un \VUgras{centinela} \textsuperscript{(5)} en su
\VUgras{garita}; parecía montar guardia sobre toda una casa de fieras de
cartón. ¡Cuántos animales había! Un \VUgras{loro} \textsuperscript{(6)} en su percha,
un \VUgras{rinoceronte} \textsuperscript{(7)} con un cuerno en la nariz, un
\VUgras{reno} \textsuperscript{(8)}, un \VUgras{avestruz} \textsuperscript{(9)}, un
\VUgras{mono} \textsuperscript{(10)} que cascaba una nuez, una \VUgras{zorra}
\textsuperscript{(11)} que se llevaba una gallina, un \VUgras{cocodrilo}
\textsuperscript{(12)} con sus enormes fauces abiertas, un \VUgras{camello}
\textsuperscript{(13)}, un elegante \VUgras{galgo} \textsuperscript{(14)}, una
\VUgras{pantera} \textsuperscript{(15)}, y luego dos bichos que yo no conocía, que
son, según me dicen, una \VUgras{comadreja} \textsuperscript{(16)} y una
\VUgras{hiena} \textsuperscript{(17)}. Había también un \VUgras{leopardo}
\textsuperscript{(18)} y un \VUgras{lobo} \textsuperscript{(21)}; muy cerca había lindas
\VUgras{ratas} \textsuperscript{(19)} y \VUgras{ratones} \textsuperscript{(20)} diminutos de
caucho. Por último, un \VUgras{hipopótamo} \textsuperscript{(22)} y un
\VUgras{dromedario} \textsuperscript{(23)} jorobado completaban la colección. Me
habría quedado allí horas enteras si no hubiera habido, muy cerca, un
espléndido campamento lleno de \VUgras{soldaditos de plomo} \textsuperscript{(29)}
que me llamó la atención.

%%K t16-tit-2 sub
\VUsaut{4.0mm}
\VUpk{10.2pt}{40}{Soldados y guerra.}
\VUsaut{2.8mm}

%%K t16-04-1 p
4. --- Mi abuelo, que es un \VUgras{inválido de guerra} \textsuperscript{(24)} y tuvo
que dejar el ejército a raíz de una \VUgras{herida} que le valió la
\VUgras{medalla militar}, adora contarme las \VUgras{campañas} en que
participó. Le encantó explicarme los distintos movimientos del pequeño
ejército en miniatura. Un grupo de soldados de verdad que visitaban el
bazar --- un \VUgras{ingeniero militar} \textsuperscript{(25)}, un
\VUgras{cazador alpino} \textsuperscript{(26)} con su \VUgras{boina}, un
\VUgras{sargento primero} \textsuperscript{(27)} de infantería y un \VUgras{sargento}
\textsuperscript{(28)} de artillería con un \VUgras{galón} de oro en la manga --- se
detuvo cerca de nosotros a escuchar las explicaciones.

%%K t16-05-1 p
5. --- Mi abuelo, muy orgulloso de tener un público tan distinguido,
improvisó todo un plan de batalla.

%%K t16-05-2 p
La plaza fuerte, con sus \VUgras{murallas} y sus \VUgras{fosos}, estaba
sitiada por el \VUgras{ejército enemigo}, que, tras un largo
\VUgras{bombardeo}, se disponía a tomarla por \VUgras{asalto}. Pero los
\VUgras{sitiados}, empujados por el hambre, habían intentado una
\VUgras{salida} contra el \VUgras{campamento} de los \VUgras{sitiadores}.
Rechazado el \VUgras{ataque} en un encarnizado \VUgras{combate} en torno a
las \VUgras{tiendas}, los sitiadores \VUgras{vencedores} lanzaron sus
\VUgras{escuadrones} en \VUgras{persecución} de los asaltantes
\VUgras{vencidos}. Éstos \VUgras{se retiraron}, dejando el
\VUgras{campo de batalla} sembrado de \VUgras{muertos} y \VUgras{heridos}.
Los \VUgras{camilleros} llevaron a los heridos al \VUgras{hospital de campaña}
para que los curara el \VUgras{cirujano}.

%%K t16-05-3 p
Pese a la heroica resistencia de unos cuantos \VUgras{batallones} que se
habían formado en \VUgras{cuadro}, la \VUgras{derrota} fue completa; el
resto del ejército \VUgras{huyó}, dejando atrás muchos \VUgras{prisioneros}.
El enemigo estaba a punto de coronar su \VUgras{victoria} tomando la plaza.
Quizá fuera aquél el fin de la campaña, y tras un \VUgras{armisticio} podría
firmarse un \VUgras{tratado de paz}.

%%K t16-tit-3 sub
\VUsaut{4.4mm}
\VUpk{10.2pt}{40}{Los regalos de año nuevo.}
\VUsaut{2.8mm}

%%K t16-06-1 p
6. --- Los jóvenes soldados, que se reían escuchando el pintoresco relato
del veterano, le hicieron algunas preguntas más. En cuanto a mí, di la
vuelta al mostrador para mirar una hermosa \VUgras{ballesta} \textsuperscript{(32)} y
un \VUgras{arco} \textsuperscript{(33)} con su \VUgras{flecha} \textsuperscript{(31)}. Allí
encontré otra colección de animales: dos \VUgras{pájaros cantores}
\textsuperscript{(34)}, un \VUgras{ruiseñor} y una \VUgras{curruca} de cuerda que
cantaban a voz en grito, y una serie de bichos extraños --- un
\VUgras{murciélago} \textsuperscript{(35)}, una \VUgras{boa} \textsuperscript{(36)}, una
\VUgras{tortuga} \textsuperscript{(37)}, una \VUgras{foca} \textsuperscript{(38)}, una
\VUgras{ballena} \textsuperscript{(39)} y un \VUgras{tiburón} \textsuperscript{(40)} hechos de
piel de batihoja.

%%K t16-07-1 p
7. --- No me detuve mucho a mirarlos, pues había visto lo que soñaba: un
magnífico \VUgras{teatro de títeres} \textsuperscript{(41)} con espléndidos
\VUgras{decorados} y un delicioso \VUgras{telón} rojo. En el
\VUgras{escenario}, dos actores parecían representar un \VUgras{papel} en
una \VUgras{obra} interesantísima, pues los pequeños \VUgras{espectadores}
de cartón de los palcos, o sentados en las \VUgras{butacas} y en el
\VUgras{gallinero} detrás del \VUgras{director de orquesta}, aplaudían con
entusiasmo.

%%K t16-07-2 p
Por desgracia, aquel teatro era carísimo.

%%K t16-08-1 p
8. --- Además, era difícil elegir entre los juguetes, pues los escaparates
estaban llenos de juguetes de todas clases: \VUgras{Arlequín}
\textsuperscript{(42)}, \VUgras{Polichinela} \textsuperscript{(43)}, \VUgras{Pierrot}
\textsuperscript{(44)}, \VUgras{lechuzas} \textsuperscript{(45)} que batían las alas,
\VUgras{mirlos} \textsuperscript{(46)} que silbaban, \VUgras{perros de lanas} que
ladraban, un \VUgras{gramófono} \textsuperscript{(47)} y \VUgras{rompecabezas}. Uno de
esos juguetes me divirtió muchísimo: representaba un
\VUgras{combate naval} \textsuperscript{(48)} en el que un \VUgras{barco} es atacado
por \VUgras{piratas} frente a una costa plantada de \VUgras{cocoteros}.

%%K t16-noto-1 noto
\VUnotes{143.2mm}{%
\textit{(*) Véase el cuadro n.º 6.}%
}

%%K t16-09-1 p
9. --- Pude mirar todas estas maravillas con toda calma, pues había poca
gente en los almacenes --- apenas un puñado de curiosos, un
\VUgras{dragón} \textsuperscript{(49)} y un \VUgras{cobrador} \textsuperscript{(50)} que
presentaba una \VUgras{letra de cambio} en la \VUgras{caja} \textsuperscript{(51)}
principal.

%%K t16-10-1 p
10. --- Pregunté a mi abuelo para qué servían los libros que había en los
estantes detrás del \VUgras{cajero}; me dijo que eran los
\VUgras{libros de cuentas}.

%%K t16-tit-4 sub
\VUsaut{3.6mm}
\VUpk{10.2pt}{40}{Mirando por toda la tienda.}
\VUsaut{2.8mm}

%%K t16-11-1 p
11. --- Al salir de la sección de juguetes, fuimos a la guarnicionería a
echar un vistazo a las \VUgras{sillas de montar} \textsuperscript{(53)}, los
\VUgras{estribos} \textsuperscript{(54)}, las \VUgras{bridas} \textsuperscript{(55)}, los
\VUgras{bocados} \textsuperscript{(56)} y las \VUgras{fustas}. Mientras el
\VUgras{jefe de sección} \textsuperscript{(57)} daba a mi abuelo algunos informes, fui
a mirar la exposición de \VUgras{porcelana} y \VUgras{cristalería}
\textsuperscript{(58)}, donde había hermosas \VUgras{ensaladeras} y delicadas
\VUgras{teteras} \textsuperscript{(59)}.

%%K t16-12-1 p
12. --- Para subir al primer piso, pasamos por detrás del escaparate de los
\VUgras{corsés} \textsuperscript{(60)}, junto a la mercería, donde había expuestos
hermosísimos \VUgras{calzones} \textsuperscript{(63)}. Cerca, un
\VUgras{dependiente} \textsuperscript{(61)} ayudaba a una \VUgras{clienta}
\textsuperscript{(62)} a escoger hermosas \VUgras{telas} \textsuperscript{(64)} de seda.

%%K t16-12-2 p
Al subir la escalera, noté que los almacenes estaban adornados con
\VUgras{farolillos} \textsuperscript{(65)} japoneses, \VUgras{sombrillas}
\textsuperscript{(66)} y enormes \VUgras{palmeras} \textsuperscript{(70)}.

%%K t16-13-1 p
13. --- En el primer piso no había mucho que ver: sólo muebles (*),
\VUgras{cajas de caudales} \textsuperscript{(67)}, \VUgras{cómodas} \textsuperscript{(68)} y
\VUgras{cochecitos de niño} \textsuperscript{(69)}. Pero la vista general de los
almacenes era muy \VUgras{entretenida}.

%%K t16-14-1 p
14. --- Debajo de nosotros veía los instrumentos de música: \VUgras{arpas}
\textsuperscript{(71)}, \VUgras{guitarras} \textsuperscript{(72)}, \VUgras{mandolinas}
\textsuperscript{(73)}, \VUgras{cornetines} \textsuperscript{(74)}, \VUgras{clarinetes}
\textsuperscript{(75)}, violines con sus \VUgras{arcos} \textsuperscript{(76)} e incluso un
\VUgras{bombo} \textsuperscript{(77)}. Un poco más allá, en el mostrador de
papelería, un \VUgras{estudiante} \textsuperscript{(78)} regateaba con el
\VUgras{dependiente} \textsuperscript{(79)} una cartera de documentos. Cerca de ellos
distinguía un hermoso \VUgras{abecedario} \textsuperscript{(80)}.

%%K t16-14-2 p
En medio de los almacenes se alzaba un alto \VUgras{árbol de Navidad}
\textsuperscript{(81)} cargado de velas, chucherías y juguetes de todas clases:
\VUgras{caballos de madera} \textsuperscript{(82)}, \VUgras{muñecas} \textsuperscript{(83)},
etc.

%%K t16-14-3 p
La \VUgras{perfumería} \textsuperscript{(84)}, con sus \VUgras{dentífricos} y sus
diversos \VUgras{perfumes}, despedía un olor muy agradable que subía hasta
nosotros.

%%K t16-tit-5 sub
\VUsaut{3.4mm}
\VUpk{10.2pt}{40}{Compramos algo.}
\VUsaut{2.8mm}

%%K t16-15-1 p
15. --- Como mi abuelo empezaba a encontrar el tiempo largo, volvimos a
bajar por la gran escalera. Se detuvo un momento delante de un
\VUgras{cuadro de astronomía} \textsuperscript{(85)} que representaba, según me dijo,
\VUgras{cometas} \textsuperscript{(a)}, \VUgras{eclipses de luna y de sol}
\textsuperscript{(b)}, una rosa de los vientos con los \VUgras{cuatro puntos
cardinales} \textsuperscript{(c)} \VUgras{(norte, sur, este y oeste)}, un mapa de los
dos \VUgras{hemisferios} \textsuperscript{(d)}, los \VUgras{planetas} \textsuperscript{(e)} y
el \VUgras{sistema solar} \textsuperscript{(f)}. No entendí nada de aquello, y me
interesaron mucho más la librea galoneada de un \VUgras{repartidor}
\textsuperscript{(86)} que pasaba, y los bonitos \VUgras{abanicos} \textsuperscript{(87)} que
tenía al lado, junto a los \VUgras{monederos} \textsuperscript{(88)} y las
\VUgras{carteras} \textsuperscript{(89)}.

%%K t16-16-1 p
16. --- Admiré también en el mostrador unas \VUgras{plumas} \textsuperscript{(90)} y
unas \VUgras{hebillas de cinturón} \textsuperscript{(91)}, y las bonitas
\VUgras{flores artificiales} \textsuperscript{(94)} graciosamente dispuestas en
cestas. Las \VUgras{violetas}, los \VUgras{alhelíes}, los \VUgras{jacintos}
\textsuperscript{(95)}, los \VUgras{geranios} \textsuperscript{(96)}, los \VUgras{lirios}
\textsuperscript{(97)} y las \VUgras{capuchinas} \textsuperscript{(98)} estaban muy bien
imitados.

%%K t16-tit-6 sub
\VUsaut{4.4mm}
\VUpk{10.2pt}{40}{El regalo del abuelo.}
\VUsaut{2.8mm}

%%K t16-17-1 p
17. --- Pedí a mi abuelo que me comprara uno de los
\VUgras{cepillos de dientes} \textsuperscript{(92)} de una caja que había junto a las
\VUgras{horquillas} \textsuperscript{(93)}. Aceptó, y fui yo mismo a pagar mi compra
a la caja, al lado de la cual estaban haciendo un gran \VUgras{paquete}
\textsuperscript{(100)} para una \VUgras{clienta} \textsuperscript{(99)}.

%%K t16-18-1 p
18. --- De pronto hubo un ligero revuelo entre la gente que estaba junto a
los \VUgras{bronces} \textsuperscript{(101)} artísticos. Un \VUgras{ladrón}
\textsuperscript{(102)} acababa de ser sorprendido en flagrante delito por un
\VUgras{inspector de los almacenes} \textsuperscript{(103)} cuando intentaba robar a
alguien. Se mandó buscar a un guardia para \VUgras{detenerlo}, y en el
momento en que salíamos, se llevaban al culpable a la comisaría.
\VUsaut{6.0mm}
\VUfilet{22mm}
